% 迁移来源：NSFC_2026_CIESM-STSR/extraTex/1.1.立项依据.tex
\subsubsection{研究背景与意义}

新疆山地—盆地相间，高程、坡向、遮蔽及下垫面差异，使气温、风、短波辐射和降水在同一粗网格内仍有系统性差异。忽略这些差异及变量共同变化，会影响资源、生态与极端气象判断。本课题重建山地—盆地过渡带和天山南北坡梯度等宏观地形尺度气象场，不涉及坡面单点条件。

\paragraph{新疆地形对气象变量的影响}
四类机制为联合重建提供依据：盆地冬季逆温使温度—高程关系随季节反转，单一递减率难以描述；天山迎风坡抬升主导降水，高程、坡向影响积雪格局\cite{Li2019TianshanSnow}，冬春积雪又通过反照率连接短波吸收与气温；绿洲冷岛使同网格绿洲与荒漠气温相差数摄氏度\cite{Hao2015Oasis}；塔里木盆地沙尘显著削减地表太阳辐射，是短波的重要区域误差源\cite{Li2020TaklimakanDust}。这些机制使新疆适于检验地形、下垫面条件对变量依赖及重建的影响。

清华大学联合地球系统模式（Community Integrated Earth System Model，CIESM）已形成模式发展与历史试验成果，但仍有系统偏差\cite{Lin2020CIESM}。第六次国际耦合模式比较计划（CMIP6）的资料标准支持追溯试验、成员、变量和版本\cite{Eyring2016CMIP6}。本课题拟使用 CIESM 的 \texttt{historical} 输出，将约 1$^{\circ}$ 大气网格增强至 0.1$^{\circ}$（线性约 10 倍、面积约 100 倍），第一季度核定成员、网格和变量频率。该处理学习既有模式输出条件下的区域细节，不能替代模式自身分辨率提高，后者的计算成本也超出课题规模。

重建须兼顾地形与变量联系所含的有效细节，以及回聚合后的大尺度一致性；纹理清晰或单变量均值匹配均不足以证明有效。图~\ref{fig:background} 概括这一问题。研究以配对实验、模式应用和独立台站观测检验信息含量、地形与共享表征增益及约束兼容性，形成可接入 CIESM 流程的模块。
\CiesmFigure{figures/F01-background.pdf}{新疆四变量跨尺度重建：粗尺度与地形条件用于重建，虚线表示向校正状态 $X^{\star}=B(X)$ 回聚合。0.1$^{\circ}$ 网格细化不等于精度提升；地形、纹理及局地机制为示意。三类验证分别提供证据，真实 CIESM 与观测不逐日强行配对}{fig:background}

课题期限为 2027 年一年，拟申请 P02“算法或数据集”方向、经费 \CiesmValue{total-budget} 万元，按资源条件聚焦可复现方法、约束效应与接口验证，由实验判断细网格的信息增量。

\subsubsection{国内外研究现状及发展动态分析}

\paragraph{参考资料与跨尺度对应关系}
欧洲中期天气预报中心第五代再分析（ERA5）通过模式与观测同化提供连续资料\cite{Hersbach2020ERA5}。ERA5 陆面再分析（ERA5-Land）网格更细，但由 ERA5 强迫驱动，二者并非独立证据，须逐变量辨别模拟、插值与修正形成的细节\cite{Munoz2021ERA5Land}。

\paragraph{ERA5-Land 的适用范围与参考场选择}
ERA5-Land 气温含高程递减率订正形成的结构；降水、风与地表向下短波主要由 ERA5 插值，缺少地形抬升、遮蔽或坡向订正\cite{Munoz2021ERA5Land}，统一作为真值可能只学到插值纹理。因此，气温采用 ERA5-Land，另三变量采用融合台站、卫星和再分析的中国区域地面气象要素驱动数据集（CMFD）\cite{He2020CMFD}，相应 ERA5-Land 变量仅作对照；中国气象数据网地面气候日值产品中的新疆国家级台站用于独立验证，不参与训练。三类资料均按公开许可获取，无需专项授权；第一季度核定年份、频率和站点可用性。

自由运行气候模拟与再分析同年不同天气，不能按日期拼接监督样本。DeepESD 以再分析学习粗细尺度关系再迁移至气候模式，体现 perfect-prognosis 思路\cite{Bano2022DeepESD}，但仍依赖预测因子分布适应性和关系稳定性。

\paragraph{从空间纹理到多变量时空重建}
DeepSD 引入卷积超分辨率\cite{DeepSD2017}，Stengel 等扩展至风和太阳辐射的对抗式增强\cite{Stengel2020}，Serifi 等揭示变量类型与采样频率影响时空重建结构选择\cite{Serifi2021}，提示平滑温度与间歇降水应采用不同误差模型。生成方法进一步刻画条件细尺度分布：Harris 等用概率评分、空间谱与校准诊断降水随机降尺度\cite{Harris2022}；CorrDiff 分解条件均值与随机残差\cite{Mardani2025CorrDiff}；Schmidt 等研究概率化、时空连续的气候降尺度\cite{Schmidt2025Coherent}。现有多变量生成成果之上，仍须实验证明本课题地形与约束的增益。

\paragraph{偏差校正、变量依赖与非配对迁移}
单变量分位数校正不能替代依赖结构评价。多变量分位数映射\cite{Cannon2018MBCn}、区分趋势保持与校正降尺度职责的 ISIMIP3BASD\cite{Lange2019ISIMIP}，以及深度超分辨率结合趋势保持校正的方法\cite{Trend2024}提供了参照；研究须纳入竞争性统计基线，并报告校正造成的模式状态变化。

非配对研究包括扩散桥变换\cite{Bischoff2024Unpaired}、粗尺度去偏后条件采样\cite{Wan2023}及带不确定性刻画的快速尺度自适应生成降尺度\cite{Hess2025Scale}。Aich 等以条件扩散联合校正、降尺度模式降水，无需原始模式与观测轨迹强行配对\cite{Aich2026Diffusion}，但尚未回答跨变量依赖恢复问题。非配对框架仍须明确保留尺度、调整统计量及区域适用性。

\paragraph{一致性约束与泛化评价}
面积加权聚合约束\cite{Geiss2022}和温度变量硬约束\cite{Gonzalez2024Hard}表明，可定义的数学条件可嵌入输出，但依赖变量定义与观测支持；本课题四变量不足以闭合完整水量、能量收支。过程合理性\cite{Maraun2017}、较暖气候降水极端外推\cite{Rampal2024}和跨参考数据集适应性\cite{Pretrain2026}研究提示：评价须覆盖连续时段、地形分层及域外区域，不能仅看随机划分的平均误差。

\subsubsection{现有研究的不足}

第一，细网格不必然含独立信息，须逐变量检验细节来源。人工降质配对样本未必覆盖 CIESM 偏差，模式应用后的分布改善也不能证明逐日过程正确，两类证据须分别解释。

第二，辅助信息、容量与约束可能混淆增益来源，须在相同数据、参数量级和计算预算下分离变量共享、地形与一致性条件的贡献，并检查降水等变量的负迁移。

第三，与原始 CIESM 强制一致可能保留偏差，与校正场一致则须报告信号变化。温度垂直梯度、相关性和空间梯度随季节、天气型变化，不能作为普适硬约束；其与有效细节的权衡须经独立对照验证。

\subsubsection{关键科学问题与研究切入点}

中心问题是：\textbf{在 CIESM 与参考资料存在分布差异时，复杂地形和多变量依赖能否提供稳定的亚网格重建增益，并与大尺度一致性条件兼容？}

具体追问：校正粗尺度状态后，哪些时空结构可学习、可迁移？地形条件下的变量共享何时在山区、盆地产生增益或失效？区分聚合条件与经验约束后，如何辨别限制非物理输出与过度压平局地变化的界限？

可证伪假设为：固定资料、校正与容量时，地形条件下的共享表征有望减少有限样本估计误差和条件重建误差，改善变量共同变化；适当聚合及变量定义约束有望在保留增益时降低不一致残差。增益定义为等信息独立基线误差减去本方法误差，正值表示改善。

唯一主要评价量为无量纲地形分层空间结构误差，其最小有意义增益（MCID）为等信息独立基线验证集指标的 5\%；主要增益为同一测试样本上基线与本方法该误差之差。95\% 区间上界低于 MCID 时不支持有效增益；下界高于 MCID 且四变量均非退化时才支持；其余为证据不足，任一变量退化越界均不判整体有效。验证集基线误差为零时，MCID 取零、单列说明且不算相对改善率。

逐变量信息含量与误差改善以平均绝对误差（MAE）衡量，MCID 取等信息独立基线验证集均方根误差（RMSE）的 5\% 与绝对下限的较大者：气温 0.15\,K、近地面风 0.20\,m$\cdot$s$^{-1}$、地表向下短波 3.0\,W$\cdot$m$^{-2}$、降水 0.30\,mm$\cdot$d$^{-1}$。RMSE 仅定尺度，增益始终为对照 MAE 减去本方法 MAE；信息含量检验对照为仅插值，联合重建对照为等信息独立模型。非退化要求各变量 MAE 增加不超过独立基线验证集 RMSE 的 3\%。地形与共享的交互不成立即撤回地形增强解释。计算规则及绝对下限在申请阶段确定，训练、验证资料所得尺度和阈值在查看测试集前固定存档，之后不再选择。

研究依次建立可追溯资料、构建联合模型，以单变量、无地形和不同约束强度对照识别增益，完成 CIESM 分布、结构与接口验证，阐明地形、变量依赖及尺度条件的适用边界，形成输入输出可检查的增强模块。

\printbibliography[heading=ciesm]
